
\documentclass[11pt,a4paper]{article}

\usepackage[utf8]{inputenc}
\usepackage[T1]{fontenc}
\usepackage[french]{babel}
\usepackage{lmodern}
\usepackage{microtype}
\usepackage[a4paper,margin=1.7cm,headheight=15pt]{geometry}
\usepackage{amsmath,amssymb,mathtools}
\usepackage{array,tabularx,multicol}
\usepackage{enumitem}
\usepackage[most]{tcolorbox}
\usepackage{xcolor}
\usepackage{tikz}
\usetikzlibrary{arrows.meta,calc}
\usepackage{fancyhdr}
\usepackage{listings}
\usepackage{hyperref}

\definecolor{bleufonce}{RGB}{28,55,120}
\definecolor{bleuclair}{RGB}{238,243,255}
\definecolor{grisclair}{RGB}{246,246,246}
\definecolor{tableline}{RGB}{45,45,45}
\definecolor{softgray}{RGB}{246,247,249}
\definecolor{deepblue}{RGB}{25,62,115}
\definecolor{softblue}{RGB}{235,241,250}
\definecolor{accent}{RGB}{20,82,130}

\hypersetup{
  colorlinks=true,
  linkcolor=bleufonce,
  urlcolor=bleufonce,
  pdftitle={Parcours de révision -- Probabilités -- Corrigés},
  pdfauthor={}
}

\setlength{\parindent}{0pt}
\setlength{\parskip}{0.28em}
\renewcommand{\arraystretch}{1.12}
\setlist[itemize]{leftmargin=1.25em,itemsep=0.15em,topsep=0.2em}
\setlist[enumerate,1]{label=\textbf{\arabic*.}, leftmargin=1.05cm, itemsep=0.35em, topsep=0.25em}
\setlist[enumerate,2]{label=\textbf{\alph*.}, leftmargin=0.85cm, itemsep=0.25em, topsep=0.2em}

\pagestyle{fancy}
\fancyhf{}
\lhead{}
\rhead{}
\cfoot{\thepage}

\newcommand{\Bin}{\mathcal B}
\newcommand{\origine}[1]{ {\scriptsize\textcolor{bleufonce}{(site #1)}}}
\newcounter{reponse}

\newenvironment{blocreponses}[1]{%
  \begin{tcolorbox}[
    enhanced,
    breakable,
    colback=bleuclair,
    colframe=bleufonce,
    boxrule=0.6pt,
    arc=2mm,
    left=2mm,right=2mm,top=2mm,bottom=1.5mm,
    title={\bfseries #1},
    coltitle=white,
    colbacktitle=bleufonce,
    before skip=3mm,
    after skip=3mm, fontupper=\large
  ]
  \large
  \begin{enumerate}[leftmargin=*,label=\textbf{\arabic*.},itemsep=0.45em,topsep=0.5em]
  \setcounter{enumi}{\value{reponse}}
}{%
  \setcounter{reponse}{\value{enumi}}
  \end{enumerate}
  \end{tcolorbox}
}

\newtcolorbox{correctionbox}[1]{
  enhanced,
  breakable,
  colback=grisclair,
  colframe=black!55,
  boxrule=0.55pt,
  arc=2mm,
  left=2mm,right=2mm,top=2mm,bottom=1.5mm,
  title=\textbf{#1},
  coltitle=black,
  colbacktitle=black!12,
  before skip=2mm,
  after skip=2mm
}

\newcommand{\titreSujet}[2]{%
  \subsection*{#1}
  \addcontentsline{toc}{subsection}{#1}
  \textit{Référence : #2}\par\medskip\hrule\medskip
}
\newcommand{\vect}[1]{\overrightarrow{#1}}
\newcommand{\e}{\mathrm e}

\tikzset{
  tabouter/.style={draw=tableline, line width=0.85pt},
  tabline/.style={draw=tableline, line width=0.55pt},
  forbidden/.style={draw=tableline, line width=0.95pt},
  vararrow/.style={-{Latex[length=2.7mm,width=2.0mm]}, line width=0.85pt, draw=accent},
  tabtext/.style={font=\small},
  tabmath/.style={font=\small},
  labelcell/.style={font=\small\bfseries, text=deepblue},
  pointvalue/.style={font=\small\bfseries},
  endvalue/.style={font=\small},
}

\lstset{language=Python,basicstyle=\ttfamily\small,numbers=left,numberstyle=\tiny,frame=single,showstringspaces=false,columns=fullflexible,keepspaces=true}

\begin{document}

\begin{center}
{\LARGE\bfseries Parcours de révision -- Probabilités}\\[1mm]

\end{center}

\vspace{1mm}
\tableofcontents
\clearpage

\section*{Partie 1 -- Corrections des questions flashs}
\addcontentsline{toc}{section}{Partie 1 -- Corrections des questions flashs}
\setcounter{reponse}{0}

\begin{blocreponses}{Réponses -- A. Vocabulaire essentiel}
  \item L'univers est l'ensemble de toutes les issues possibles.
  \item Une expérience aléatoire est une expérience dont on ne peut pas prévoir avec certitude l'issue.
  \item Une issue est un résultat possible d'une expérience aléatoire.
  \item La loi de probabilité d'une variable aléatoire donne les valeurs prises par cette variable et les probabilités associées.
  \item Une variable aléatoire réelle est une grandeur numérique dépendant du hasard.
  \item L'évènement contraire de $A$ se note souvent $\overline A$.
  \item $A\cap B$ signifie : « $A$ et $B$ ».
  \item $A\cup B$ signifie : « $A$ ou $B$ ».
  \item Deux évènements incompatibles ne peuvent pas se produire en même temps : $A\cap B=\varnothing$.
  \item Deux évènements sont indépendants lorsque $P(A\cap B)=P(A)P(B)$.
  \item Une partition de l'univers est une famille d'évènements incompatibles deux à deux dont la réunion est l'univers.
  \item Une épreuve de Bernoulli est une expérience aléatoire à deux issues : succès ou échec.
\end{blocreponses}

\begin{blocreponses}{Réponses -- B. Premiers calculs et arbres de probabilités}
  \item $P(\overline A)=1-P(A)$.
  \item La probabilité de l'évènement impossible est $0$.
  \item La probabilité de l'évènement certain est $1$.
  \item $P(\overline A)=1-0{,}35=0{,}65$.
  \item Si $A$ et $B$ forment une partition de l'univers, alors $P(A)+P(B)=1$.
  \item La probabilité de l'échec est $1-0{,}2=0{,}8$.
  \item Si $P(B)\neq0$, alors $P_B(A)=\dfrac{P(A\cap B)}{P(B)}$.
  \item $P(A\cap B)=P(B)\times P_B(A)$, ou encore $P(A\cap B)=P(A)\times P_A(B)$.
  \item $P_B(A)=\dfrac{P(A\cap B)}{P(B)}=\dfrac{\frac15}{\frac3{10}}=\dfrac15\times\dfrac{10}{3}=\dfrac23$.
  \item On calcule la probabilité d'un chemin en multipliant les probabilités inscrites sur ses branches.
  \item On additionne les probabilités des différents chemins qui réalisent l'évènement.
  \item Si les évènements $B_i$ forment une partition de l'univers, alors $P(A)=\displaystyle\sum_i P(B_i\cap A)=\sum_i P(B_i)P_{B_i}(A)$.
\end{blocreponses}

\begin{blocreponses}{Réponses -- C. Loi binomiale}
  \item Un schéma de Bernoulli est une répétition d'épreuves de Bernoulli identiques et indépendantes, avec la même probabilité de succès.
  \item Une variable aléatoire suit une loi binomiale lorsqu'elle compte le nombre de succès dans un schéma de Bernoulli.
  \item On note $X\sim\mathcal B(n,p)$.
  \item Une variable $X\sim\mathcal B(n,p)$ peut prendre toutes les valeurs entières de $0$ à $n$.
  \item Elle compte le nombre de succès.
  \item $P(X=k)=\binom{n}{k}p^k(1-p)^{n-k}$.
  \item $P(X=3)=\binom83\left(\dfrac14\right)^3\left(\dfrac34\right)^5$.
  \item Pour $X\sim\mathcal B(n,p)$, $E(X)=np$.
  \item Dans un contexte concret, l'espérance représente une valeur moyenne attendue sur un grand nombre de répétitions.
  \item $E(X)=10\times0{,}2=2$.
  \item Pour $X\sim\mathcal B(n,p)$, $V(X)=np(1-p)$.
  \item $\sigma(X)=\sqrt{V(X)}=\sqrt{np(1-p)}$.
  \item On utilise souvent $P(X\geqslant1)=1-P(X=0)$.
\end{blocreponses}

\begin{blocreponses}{Réponses -- D. Variables aléatoires finies}
  \item Si $X$ prend les valeurs $x_1,\ldots,x_n$, alors $E(X)=\displaystyle\sum_{i=1}^n x_iP(X=x_i)$.
  \item $V(X)=E\bigl((X-E(X))^2\bigr)$, ou $V(X)=\displaystyle\sum_i (x_i-E(X))^2P(X=x_i)$.
  \item $\sigma(X)=\sqrt{V(X)}$.
  \item $E(X)=0\times\dfrac14+1\times\dfrac12+3\times\dfrac14=\dfrac12+\dfrac34=\dfrac54=1{,}25$.
  \item Les probabilités sont positives et leur somme vaut $1$.
  \item On note cet évènement $\{X=a\}$.
  \item On additionne les probabilités des valeurs de $X$ inférieures ou égales à $a$.
\end{blocreponses}

\begin{blocreponses}{Réponses -- E. Rédaction type bac}
  \item Il faut préciser les épreuves de Bernoulli, l'indépendance, le caractère identique des répétitions, le succès et sa probabilité, puis dire que $X$ compte le nombre de succès. On conclut alors $X\sim\mathcal B(n,p)$.
  \item Il faut annoncer la partition avant la formule. Par exemple : « $A$ et $\overline A$ forment une partition de l'univers. D'après la formule des probabilités totales, on a : $P(B)=P(A\cap B)+P(\overline A\cap B)$. »
  \item On peut écrire : « On répète $10$ expériences de Bernoulli identiques et indépendantes, de même probabilité de succès $p$. La variable aléatoire $X$ compte le nombre de succès obtenus. Donc $X\sim\mathcal B(10,p)$. »
  \item Il faut annoncer la partition avant la formule : « $A$ et $\overline A$ forment une partition de l'univers. D'après la formule des probabilités totales, $P(B)=P(A\cap B)+P(\overline A\cap B)$. »
  \item Il faut ajouter l'hypothèse : « Si $A$ et $B$ sont indépendants, alors $P(A\cap B)=P(A)P(B)$. » Sans indépendance, cette égalité n'est pas vraie en général.
  \item L'évènement contraire de « au moins un succès » est « aucun succès ». Donc $P(X\geqslant1)=1-P(X=0)$.
\end{blocreponses}

\clearpage
\section*{Partie 2 -- Corrigés des sujets de bac}
\addcontentsline{toc}{section}{Partie 2 -- Corrigés des sujets de bac}

\titreSujet{Sujet 1 -- Groupes sanguins et donneurs universels}{Métropole, juin 2025, sujet 1, exercice 1}

\begin{correctionbox}{1. Arbre pondéré}
On complète d'abord les probabilités des quatre groupes sanguins :
\[
p(O)=1-p(A)-p(B)-p(AB)=1-0{,}45-0{,}10-0{,}03=0{,}42.
\]
Les probabilités conditionnelles contraires valent :
\[
p_A(\overline R)=0{,}15,\qquad p_B(\overline R)=0{,}16,\qquad p_{AB}(\overline R)=0{,}18.
\]

\begin{center}
\begin{tikzpicture}[grow=right, level distance=2.8cm, sibling distance=18mm,
  every node/.style={font=\small}, edge from parent/.style={draw},
  edge from parent path={(\tikzparentnode.east) -- (\tikzchildnode.west)}]
\node {}
  child {node {$O$}
    child {node {$\overline R$}}
    child {node {$R$}}
    edge from parent node[below] {$0{,}42$}}
  child {node {$AB$}
    child {node {$\overline R$} edge from parent node[below] {$0{,}18$}}
    child {node {$R$} edge from parent node[above] {$0{,}82$}}
    edge from parent node[below] {$0{,}03$}}
  child {node {$B$}
    child {node {$\overline R$} edge from parent node[below] {$0{,}16$}}
    child {node {$R$} edge from parent node[above] {$0{,}84$}}
    edge from parent node[above] {$0{,}10$}}
  child {node {$A$}
    child {node {$\overline R$} edge from parent node[below] {$0{,}15$}}
    child {node {$R$} edge from parent node[above] {$0{,}85$}}
    edge from parent node[above] {$0{,}45$}};
\end{tikzpicture}
\end{center}
\end{correctionbox}

\begin{correctionbox}{2. Calcul de $p(B\cap R)$}
On utilise la formule d'une probabilité d'intersection avec une probabilité conditionnelle :
\[
p(B\cap R)=p(B)\times p_B(R)=0{,}10\times 0{,}84=0{,}084.
\]
Ainsi, la probabilité qu'une personne choisie au hasard soit de groupe sanguin B et de rhésus positif est égale à $0{,}084$, soit $8{,}4\,\%$.
\end{correctionbox}

\begin{correctionbox}{3. Calcul de $p_O(R)$}
Les événements $A$, $B$, $AB$ et $O$ forment une partition de l'univers. D'après la formule des probabilités totales, on a :
\[
p(R)=p(A\cap R)+p(B\cap R)+p(AB\cap R)+p(O\cap R).
\]
Donc
\[
0{,}8397=0{,}45\times0{,}85+0{,}10\times0{,}84+0{,}03\times0{,}82+0{,}42\times p_O(R).
\]
Or
\[
0{,}45\times0{,}85+0{,}10\times0{,}84+0{,}03\times0{,}82=0{,}4911.
\]
Ainsi,
\[
0{,}42\times p_O(R)=0{,}8397-0{,}4911=0{,}3486,
\]
puis
\[
p_O(R)=\dfrac{0{,}3486}{0{,}42}=0{,}83.
\]
\end{correctionbox}

\begin{correctionbox}{4. Probabilité d'être donneur universel}
Un donneur universel est une personne de groupe O et de rhésus négatif. On cherche donc $p(O\cap \overline R)$.

D'après la question précédente, $p_O(R)=0{,}83$, donc $p_O(\overline R)=0{,}17$.

Ainsi,
\[
p(O\cap \overline R)=p(O)\times p_O(\overline R)=0{,}42\times0{,}17=0{,}0714.
\]
La probabilité qu'un individu choisi au hasard dans la population française soit donneur universel est donc $0{,}0714$, soit $7{,}14\,\%$.
\end{correctionbox}

\begin{correctionbox}{5. Loi binomiale}
\begin{enumerate}
  \item On répète $100$ expériences de Bernoulli indépendantes et identiques, car le choix de l'échantillon est assimilé à un tirage avec remise. Le succès est : « la personne choisie est donneur universel », de probabilité $p=0{,}0714$. La variable aléatoire $X$ compte le nombre de succès parmi les 100 personnes. Donc
  \[
  X\sim \mathcal B(100;0{,}0714).
  \]

  \item On cherche $P(X\leqslant 7)$. À la calculatrice, avec $X\sim\mathcal B(100;0{,}0714)$,
  \[
  P(X\leqslant 7)\approx 0{,}577.
  \]
  La probabilité qu'il y ait au plus 7 donneurs universels dans cet échantillon est donc environ $0{,}577$.

  \item Pour une loi binomiale $\mathcal B(n;p)$, on a $E(X)=np$ et $V(X)=np(1-p)$. Donc
  \[
  E(X)=100\times0{,}0714=7{,}14,
  \]
  et
  \[
  V(X)=100\times0{,}0714\times(1-0{,}0714)=6{,}630204\approx 6{,}63.
  \]
\end{enumerate}
\end{correctionbox}

\begin{correctionbox}{6. Moyenne sur $N$ villes et Bienaymé-Tchebychev}
\begin{enumerate}
  \item La variable aléatoire $M_N$ représente le nombre moyen de donneurs universels dans les échantillons de 100 personnes prélevés dans les $N$ villes.

  \item Par linéarité de l'espérance,
  \[
  E(M_N)=E\left(\dfrac{X_1+X_2+\cdots+X_N}{N}\right)
  =\dfrac{E(X_1)+E(X_2)+\cdots+E(X_N)}{N}.
  \]
  Comme chaque variable $X_i$ a pour espérance $7{,}14$, on obtient
  \[
  E(M_N)=\dfrac{N\times7{,}14}{N}=7{,}14.
  \]

  \item Comme les variables $X_1,\ldots,X_N$ sont indépendantes,
  \[
  V(X_1+\cdots+X_N)=V(X_1)+\cdots+V(X_N)=N\times6{,}63.
  \]
  Donc
  \[
  V(M_N)=V\left(\dfrac{X_1+\cdots+X_N}{N}\right)=\dfrac{1}{N^2}V(X_1+\cdots+X_N)
  =\dfrac{N\times6{,}63}{N^2}=\dfrac{6{,}63}{N}.
  \]

  \item On remarque que
  \[
  7<M_N<7{,}28 \quad \Longleftrightarrow \quad |M_N-7{,}14|<0{,}14.
  \]
  D'après l'inégalité de Bienaymé-Tchebychev,
  \[
  P\left(|M_N-E(M_N)|\geqslant 0{,}14\right)\leqslant \dfrac{V(M_N)}{0{,}14^2}.
  \]
  Donc
  \[
  P\left(|M_N-7{,}14|\geqslant 0{,}14\right)\leqslant \dfrac{6{,}63}{N\times0{,}0196}.
  \]
  Ainsi,
  \[
  P(7<M_N<7{,}28)\geqslant 1-\dfrac{6{,}63}{N\times0{,}0196}.
  \]
  On veut que ce minorant soit strictement supérieur à $0{,}95$ :
  \[
  1-\dfrac{6{,}63}{N\times0{,}0196}>0{,}95.
  \]
  Cela équivaut à
  \[
  \dfrac{6{,}63}{N\times0{,}0196}<0{,}05,
  \]
  donc
  \[
  N>\dfrac{6{,}63}{0{,}05\times0{,}0196}\approx6765{,}31.
  \]
  La plus petite valeur entière possible est donc
  \[
  \boxed{N=6766}.
  \]
\end{enumerate}
\end{correctionbox}

\newpage

\titreSujet{Sujet 2 -- Chutes en roller et temps d'attente}{Métropole, juin 2025, sujet 2, exercice 1}

\subsubsection*{Partie A}

\begin{correctionbox}{1. Arbre pondéré}
On a $P(A)=0{,}6$, donc $P(\overline A)=0{,}4$. De plus,
\[
P_A(B)=0{,}3,\quad P_A(\overline B)=0{,}7,
\quad P_{\overline A}(B)=0{,}4,
\quad P_{\overline A}(\overline B)=0{,}6.
\]

\begin{center}
\begin{tikzpicture}[grow=right,
  level 1/.style={level distance=2.7cm, sibling distance=30mm},
  level 2/.style={level distance=2.7cm, sibling distance=12mm},
  every node/.style={font=\small}, edge from parent/.style={draw},
  edge from parent path={(\tikzparentnode.east) -- (\tikzchildnode.west)}]
\node {}
  child {node {$\overline A$}
    child {node {$\overline B$} edge from parent node[below] {$0{,}6$}}
    child {node {$B$} edge from parent node[above] {$0{,}4$}}
    edge from parent node[below] {$0{,}4$}}
  child {node {$A$}
    child {node {$\overline B$} edge from parent node[below] {$0{,}7$}}
    child {node {$B$} edge from parent node[above] {$0{,}3$}}
    edge from parent node[above] {$0{,}6$}};
\end{tikzpicture}
\end{center}
\end{correctionbox}

\begin{correctionbox}{2. Calcul de $P(\overline A\cap\overline B)$}
\[
P(\overline A\cap\overline B)=P(\overline A)\times P_{\overline A}(\overline B)=0{,}4\times0{,}6=0{,}24.
\]
La probabilité qu'une personne ne chute ni pendant la première séance ni pendant la deuxième séance est donc égale à $0{,}24$.
\end{correctionbox}

\begin{correctionbox}{3. Calcul de $P(B)$}
Les événements $A$ et $\overline A$ forment une partition de l'univers. D'après la formule des probabilités totales, on a :
\[
P(B)=P(A\cap B)+P(\overline A\cap B).
\]
Donc
\[
P(B)=P(A)P_A(B)+P(\overline A)P_{\overline A}(B)=0{,}6\times0{,}3+0{,}4\times0{,}4.
\]
Ainsi,
\[
P(B)=0{,}18+0{,}16=0{,}34.
\]
\end{correctionbox}

\begin{correctionbox}{4. Probabilité conditionnelle inverse}
On cherche $P_{\overline B}(\overline A)$.

D'après la formule des probabilités conditionnelles,
\[
P_{\overline B}(\overline A)=\dfrac{P(\overline A\cap\overline B)}{P(\overline B)}.
\]
Or $P(\overline A\cap\overline B)=0{,}24$ et
\[
P(\overline B)=1-P(B)=1-0{,}34=0{,}66.
\]
Donc
\[
P_{\overline B}(\overline A)=\dfrac{0{,}24}{0{,}66}\approx0{,}364.
\]
La probabilité que la personne n'ait pas chuté lors de la première séance sachant qu'elle n'a pas chuté pendant la deuxième est donc environ $0{,}364$.
\end{correctionbox}

\begin{correctionbox}{5. Loi binomiale}
\begin{enumerate}
  \item On répète $100$ expériences de Bernoulli indépendantes et identiques, car le choix d'un échantillon de 100 personnes est assimilé à un tirage avec remise. Le succès est : « la personne ne chute ni pendant la première séance ni pendant la deuxième », de probabilité $p=0{,}24$. La variable aléatoire $X$ compte le nombre de succès parmi les 100 personnes. Donc
  \[
  X\sim\mathcal B(100;0{,}24).
  \]

  \item On cherche
  \[
  P(X\geqslant20)=1-P(X\leqslant19).
  \]
  À la calculatrice, avec $X\sim\mathcal B(100;0{,}24)$,
  \[
  P(X\geqslant20)\approx0{,}855.
  \]
  La probabilité d'avoir au moins 20 personnes qui ne chutent ni lors de la première ni lors de la deuxième séance est donc environ $0{,}855$.

  \item Pour une loi binomiale $\mathcal B(n;p)$, on a $E(X)=np$. Donc
  \[
  E(X)=100\times0{,}24=24.
  \]
  Sur un grand nombre d'échantillons de 100 personnes, on peut s'attendre en moyenne à environ 24 personnes qui ne chutent ni pendant la première séance ni pendant la deuxième.
\end{enumerate}
\end{correctionbox}

\subsubsection*{Partie B}

\begin{correctionbox}{1. Espérance de $T$}
Comme $T=T_1+T_2$, par linéarité de l'espérance,
\[
E(T)=E(T_1)+E(T_2)=40+60=100.
\]
Le temps total d'attente moyen pendant le week-end est donc de 100 minutes.
\end{correctionbox}

\begin{correctionbox}{2. Variance de $T$}
On sait que $\sigma(T_1)=10$, donc $V(T_1)=10^2=100$, et que $\sigma(T_2)=16$, donc $V(T_2)=16^2=256$.

Les variables aléatoires $T_1$ et $T_2$ sont indépendantes, donc
\[
V(T)=V(T_1+T_2)=V(T_1)+V(T_2)=100+256=356.
\]
\end{correctionbox}

\begin{correctionbox}{3. Inégalité de Bienaymé-Tchebychev}
On remarque que
\[
60<T<140 \quad \Longleftrightarrow \quad |T-100|<40.
\]
D'après l'inégalité de Bienaymé-Tchebychev,
\[
P(|T-E(T)|\geqslant40)\leqslant \dfrac{V(T)}{40^2}.
\]
Comme $E(T)=100$ et $V(T)=356$, on obtient
\[
P(|T-100|\geqslant40)\leqslant \dfrac{356}{1600}=0{,}2225.
\]
Donc
\[
P(|T-100|<40)\geqslant1-0{,}2225=0{,}7775.
\]
Ainsi,
\[
P(60<T<140)\geqslant0{,}7775>0{,}77.
\]
On a bien montré que la probabilité demandée est supérieure à $0{,}77$.
\end{correctionbox}

\newpage

\titreSujet{Sujet 3 -- Jeu vidéo et suite probabiliste}{Asie, juin 2024, sujet 2, exercice 2}

\begin{correctionbox}{1. Probabilité conditionnelle}
Si Léa vient de gagner une partie, elle gagne la suivante dans 70 \% des cas. Donc, sachant qu'elle a gagné la première partie, elle perd la deuxième avec probabilité
\[
p_{G_1}(D_2)=1-0{,}7=0{,}3.
\]
\end{correctionbox}

\begin{correctionbox}{2. Arbre des deux premières parties}
Comme Léa pense avoir autant de chances de gagner la première partie que de la perdre, on a $P(G_1)=0{,}5$ et $P(D_1)=0{,}5$.

\begin{center}
\begin{tikzpicture}[grow=right,
  level 1/.style={level distance=2.5cm, sibling distance=30mm},
  level 2/.style={level distance=2.5cm, sibling distance=12mm},
  every node/.style={font=\small}, edge from parent/.style={draw},
  edge from parent path={(\tikzparentnode.east) -- (\tikzchildnode.west)}]
\node {}
  child {node {$D_1$}
    child {node {$D_2$} edge from parent node[below] {$0{,}8$}}
    child {node {$G_2$} edge from parent node[above] {$0{,}2$}}
    edge from parent node[below] {$0{,}5$}}
  child {node {$G_1$}
    child {node {$D_2$} edge from parent node[below] {$0{,}3$}}
    child {node {$G_2$} edge from parent node[above] {$0{,}7$}}
    edge from parent node[above] {$0{,}5$}};
\end{tikzpicture}
\end{center}
\end{correctionbox}

\begin{correctionbox}{3. Calcul de $g_2$}
Les événements $G_1$ et $D_1$ forment une partition de l'univers. D'après la formule des probabilités totales, on a :
\[
g_2=P(G_2)=P(G_1\cap G_2)+P(D_1\cap G_2).
\]
Donc
\[
g_2=P(G_1)p_{G_1}(G_2)+P(D_1)p_{D_1}(G_2)=0{,}5\times0{,}7+0{,}5\times0{,}2.
\]
Ainsi,
\[
g_2=0{,}35+0{,}10=0{,}45.
\]
\end{correctionbox}

\begin{correctionbox}{4. Relation de récurrence}
\begin{enumerate}
  \item L'arbre général est le suivant.
  \begin{center}
  \begin{tikzpicture}[grow=right,
    level 1/.style={level distance=2.7cm, sibling distance=30mm},
    level 2/.style={level distance=2.7cm, sibling distance=12mm},
    every node/.style={font=\small}, edge from parent/.style={draw},
    edge from parent path={(\tikzparentnode.east) -- (\tikzchildnode.west)}]
  \node {}
    child {node {$D_n$}
      child {node {$D_{n+1}$} edge from parent node[below] {$0{,}8$}}
      child {node {$G_{n+1}$} edge from parent node[above] {$0{,}2$}}
      edge from parent node[below] {$1-g_n$}}
    child {node {$G_n$}
      child {node {$D_{n+1}$} edge from parent node[below] {$0{,}3$}}
      child {node {$G_{n+1}$} edge from parent node[above] {$0{,}7$}}
      edge from parent node[above] {$g_n$}};
  \end{tikzpicture}
  \end{center}

  \item Les événements $G_n$ et $D_n$ forment une partition de l'univers. D'après la formule des probabilités totales, on a :
  \[
  g_{n+1}=P(G_{n+1})=P(G_n\cap G_{n+1})+P(D_n\cap G_{n+1}).
  \]
  Donc
  \[
  g_{n+1}=g_n\times0{,}7+(1-g_n)\times0{,}2.
  \]
  Ainsi,
  \[
  g_{n+1}=0{,}7g_n+0{,}2-0{,}2g_n=0{,}5g_n+0{,}2.
  \]
\end{enumerate}
\end{correctionbox}

\begin{correctionbox}{5. Suite géométrique auxiliaire}
\begin{enumerate}
  \item Pour tout entier naturel $n$ non nul,
  \[
  v_{n+1}=g_{n+1}-0{,}4.
  \]
  Or $g_{n+1}=0{,}5g_n+0{,}2$, donc
  \[
  v_{n+1}=0{,}5g_n+0{,}2-0{,}4=0{,}5g_n-0{,}2=0{,}5(g_n-0{,}4)=0{,}5v_n.
  \]
  La suite $(v_n)$ est donc géométrique de raison $0{,}5$.

  Son premier terme est
  \[
  v_1=g_1-0{,}4=0{,}5-0{,}4=0{,}1.
  \]

  \item Comme $(v_n)$ est géométrique de premier terme $v_1=0{,}1$ et de raison $0{,}5$, on a, pour tout entier naturel $n$ non nul,
  \[
  v_n=0{,}1\times0{,}5^{n-1}.
  \]
  Or $v_n=g_n-0{,}4$, donc
  \[
  g_n=0{,}1\times0{,}5^{n-1}+0{,}4.
  \]
\end{enumerate}
\end{correctionbox}

\begin{correctionbox}{6. Variations de $(g_n)$}
Pour tout entier naturel $n$ non nul,
\[
g_{n+1}-g_n=(0{,}5g_n+0{,}2)-g_n=-0{,}5g_n+0{,}2=-0{,}5(g_n-0{,}4).
\]
Or
\[
g_n-0{,}4=0{,}1\times0{,}5^{n-1}>0.
\]
Donc
\[
g_{n+1}-g_n<0.
\]
La suite $(g_n)$ est donc strictement décroissante.
\end{correctionbox}

\begin{correctionbox}{7. Limite de $(g_n)$}
Comme $0<0{,}5<1$, on a
\[
\lim_{n\to+\infty}0{,}5^{n-1}=0.
\]
Donc
\[
\lim_{n\to+\infty}g_n=0{,}1\times0+0{,}4=0{,}4.
\]
À long terme, selon le modèle de Léa, la probabilité qu'elle gagne la $n$-ième partie de la journée se stabilise vers $0{,}4$, c'est-à-dire $40\,\%$.
\end{correctionbox}

\begin{correctionbox}{8. Seuil}
On cherche le plus petit entier $n$ tel que
\[
g_n-0{,}4\leqslant0{,}001.
\]
Or
\[
g_n-0{,}4=0{,}1\times0{,}5^{n-1}.
\]
On résout donc
\[
0{,}1\times0{,}5^{n-1}\leqslant0{,}001.
\]
Cela équivaut à
\[
0{,}5^{n-1}\leqslant0{,}01.
\]
En utilisant les logarithmes :
\[
(n-1)\ln(0{,}5)\leqslant \ln(0{,}01).
\]
Comme $\ln(0{,}5)<0$, le sens de l'inégalité change lorsqu'on divise :
\[
n-1\geqslant \dfrac{\ln(0{,}01)}{\ln(0{,}5)}\approx6{,}64.
\]
Donc $n\geqslant7{,}64$, et le plus petit entier convenable est
\[
\boxed{n=8}.
\]
Vérification :
\[
g_8-0{,}4=0{,}1\times0{,}5^7=0{,}00078125\leqslant0{,}001.
\]
\end{correctionbox}

\begin{correctionbox}{9. Fonction Python}
La condition de boucle doit être vraie tant que $g$ est strictement supérieur à $0{,}4+e$. À chaque passage dans la boucle, on calcule le terme suivant et on augmente le rang de 1.

\begin{lstlisting}
def seuil(e):
    g = 0.5
    n = 1
    while g > 0.4 + e:
        g = 0.5 * g + 0.2
        n = n + 1
    return n
\end{lstlisting}
\end{correctionbox}




\end{document}
